\section{Структура проекту, логіка завдань та вихідний код програми}

На відміну від процедурного підходу, кожне завдання лабораторної роботи було інкапсульовано 
в окремий автономний клас-компонент, що дозволило уникнути використання статичних методів у стилі мови C.

\begin{enumerate}
    \item \textbf{Перевірка впорядкованості масиву (Клас \texttt{ArrayChecker}):} 
    \\Алгоритм виконує один лінійний прохід по масиву починаючи з індексу \texttt{1}. 
    Якщо поточний елемент виявляється меншим за попередній, метод негайно повертає \texttt{false}, 
    що оптимізує швидкість виконання.
    \item \textbf{Матричне виведення FizzBuzz (Клас \texttt{FizzBuzzGame}):} 
    \\Для створення компактного та зручного для звітності інтерфейсу, логіку перевірки 
    відповідно до заданої кратності чисел (на 3, 5 та їхній добуток) було об'єднано з 
    формативаним виведенням \texttt{System.out.printf("\%-10s")}. Це дозволило згенерувати 
    результати у вигляді математичної сітки розміром 10 на 10 стовпців.
    \item \textbf{Пошук балансу сум (Клас \texttt{BalanceFinder}):} 
    \\Спочатку обчислюється повна сума 
    масиву, а під час другого лінійного проходу накопичується ліва частина. Рівновагу 
    знайдено в момент, коли ліва сума стає рівною рівно половині загальної суми.
    \item \textbf{Координація та точка входу (Клас \texttt{Main}):} 
    \\Клас містить обов'язковий для виконання Java-додатків метод \texttt{public static void main(String[] args)}. 
    Назначенням цього компонента є повне управління життєвим циклом програми ініціалізує тестові масиви даних, 
    виконує обчислення та створює екземпляри класів \texttt{ArrayChecker}, \texttt{FizzBuzzGame} \\
    та \texttt{BalanceFinder}, передає їм аргументи через конструктори та виводить фінальні результати обчислень у термінал середовища розробки.
\end{enumerate}

